\documentclass[11pt]{article}
\usepackage{aiguide}
\usepackage{textcomp}
\usepackage{listings}
\lstdefinestyle{lean}{basicstyle=\ttfamily\small,keywordstyle=\color{prussian}\bfseries,
  commentstyle=\color{muted}\itshape,columns=fullflexible,keepspaces=true,
  morekeywords={theorem,example,by,with,induction,rfl,rw,exact,fun},
  backgroundcolor=\color{paper},frame=leftline,framerule=2pt,rulecolor=\color{prussian},
  xleftmargin=8pt,framexleftmargin=6pt,aboveskip=0.6em,belowskip=0.6em}
\lstset{style=lean,upquote=true}
\guidekicker{Teaching guide}
\begin{document}
\makeguidetitle{Getting Started with Lean}{What a proof assistant is, a first proof checked by a computer, and where to go next.}

\section{What Lean is}
A \textbf{proof assistant} is a program that checks proofs. You write each step
in a precise language, and the computer accepts the proof only when every step
follows from the previous ones. Nothing is left as ``clearly'' or ``similarly''.

\textbf{Lean 4} is an open-source programming language and proof assistant.
\textbf{Mathlib} is its community library of formalized mathematics, from
algebra and number theory to analysis and topology. When you formalize a proof,
you find out exactly what it needs: which hypothesis is used where, which case
you forgot, which step you only believed.

\section{A first proof: $0 + n = n$}
In Lean the natural numbers are defined by the Peano axioms: every natural
number is either $0$ or the successor $k + 1$ of another one, and addition is
defined by recursion on the second argument, so $m + (k+1) = (m + k) + 1$. The
statement $0 + n = n$ is therefore not immediate; it needs induction on $n$.

\begin{lstlisting}
theorem my_zero_add (n : Nat) : 0 + n = n := by
  induction n with
  | zero => rfl
  | succ k ih => rw [Nat.add_succ, ih]
\end{lstlisting}

Read it line by line:
\begin{itemize}
  \item \texttt{induction n with} splits the proof into the two cases of the
    induction principle.
  \item \texttt{| zero => rfl}: the base case $0 + 0 = 0$ holds by the
    definition of addition, so \texttt{rfl} (``both sides are the same by
    definition'') closes it.
  \item \texttt{| succ k ih => \dots}: we must prove $0 + (k+1) = k + 1$, and
    \texttt{ih} is the induction hypothesis $0 + k = k$. The rewrite
    \texttt{Nat.add\_succ} turns the left side into $(0 + k) + 1$, and
    rewriting with \texttt{ih} turns it into $k + 1$. The goal is closed.
\end{itemize}

\textbf{When a step is wrong, Lean says so.} Replace the last line by
\texttt{| succ k ih => rfl} and Lean rejects the proof:

\begin{lstlisting}[keywordstyle=\color{ink}]
error: Tactic `rfl` failed: The left-hand side
  0 + (k + 1)
is not definitionally equal to the right-hand side
  k + 1
\end{lstlisting}

The error points at the exact step that does not follow. That is the whole idea:
you can never fool the checker, so a proof it accepts has no gaps.

{\small\color{muted}Both versions were checked with Lean 4.35.0 (core library
only, no Mathlib) in November 2026. Lean already contains this fact as
\texttt{Nat.zero\_add}; we prove our own copy to see how it works. The Natural
Number Game builds the same result step by step, in its own simplified style.}

\section{A learning path}
\begin{enumerate}[label=\textbf{\arabic*.}]
  \item \textbf{Play the Natural Number Game} in the browser. It rebuilds the
    natural numbers from the Peano axioms and teaches the basic tactics. It is
    available in Italian too.
  \item \textbf{Try more games} on the Lean Game Server, for example Scribble,
    which covers first-year university topics, or the Set Theory Game.
  \item \textbf{Experiment in Lean 4 Web}, the official web editor, which has
    Mathlib: paste the proof above and change it.
  \item \textbf{Install Lean} on your computer (VS Code with the Lean 4
    extension) when you want to work on longer files.
  \item \textbf{Work through Mathematics in Lean}, the standard tutorial for
    doing mathematics with Mathlib.
  \item \textbf{Go deeper} with \emph{Theorem Proving in Lean 4} (how Lean
    works underneath) or Kevin Buzzard's \emph{Formalising Mathematics}
    handbook (written for final-year undergraduates).
  \item \textbf{When you are stuck,} search Mathlib with Loogle, and ask in the
    new members channel of the Lean Zulip chat.
\end{enumerate}

\section{Links}
\newcolumntype{L}{>{\raggedright\arraybackslash}X}
\small
\begin{tabularx}{\linewidth}{@{}p{0.36\linewidth}L@{}}
Natural Number Game & \href{https://adam.math.hhu.de/\#/g/leanprover-community/nng4}{adam.math.hhu.de/\#/g/leanprover-community/nng4} \\
Lean Game Server & \href{https://adam.math.hhu.de/}{adam.math.hhu.de} \\
Lean 4 Web & \href{https://live.lean-lang.org/}{live.lean-lang.org} \\
Install Lean & \href{https://lean-lang.org/install/}{lean-lang.org/install} \\
Mathematics in Lean & \href{https://leanprover-community.github.io/mathematics_in_lean/}{leanprover-community.github.io/mathematics\_in\_lean} \\
Theorem Proving in Lean 4 & \href{https://lean-lang.org/theorem_proving_in_lean4/}{lean-lang.org/theorem\_proving\_in\_lean4} \\
Formalising Mathematics (K.~Buzzard) & \href{https://www.ma.imperial.ac.uk/~buzzard/xena/formalising-mathematics-2024/}{ma.imperial.ac.uk/\textasciitilde buzzard/xena/formalising-mathematics-2024} \\
Mathlib and its documentation & \href{https://mathlib.org/}{mathlib.org}, \href{https://leanprover-community.github.io/mathlib4_docs/}{leanprover-community.github.io/mathlib4\_docs} \\
Loogle (Mathlib search) & \href{https://loogle.lean-lang.org/}{loogle.lean-lang.org} \\
Lean Zulip chat & \href{https://leanprover.zulipchat.com/}{leanprover.zulipchat.com} \\
\end{tabularx}

\vspace{0.4em}
{\color{muted}Last checked: November 2026.}

\end{document}
